% EDITING: type inside the final {...} of each field. Change the shown font size
% (for example, font=10pt) in that same field block when needed.

\GrantSection{6}{Statistical Design, Safety Monitoring, Recruitment, and Operations}

\GrantGuidance{The NIH clinical-trial review criteria emphasize justified endpoints, interpretable design, adequate power, efficient conduct, ethical considerations, recruitment and retention, randomization or masking when appropriate, data management, safety monitoring, and feasible timelines.}

\GrantTextArea[
  id=ops-design-analysis, font=10pt, height=2.3in, maxchars=1700
]{Study design and analysis framework, including estimands and handling of intercurrent events}{The study is designed as a Phase 1, first-in-human, prospective, open-label, single-arm interventional trial with sequential safety review in adults with resectable or borderline-resectable KRAS G12-mutated PDAC \citep{phase1ind,onpremwhippl}. The primary analysis is estimand-based and focuses on the risk of the prespecified composite safety/feasibility endpoint under the treatment strategy of perioperative daraxonrasib plus the governed robotic procedure. Intercurrent events will be handled according to^^Jclinical relevance rather than convenience: dose reductions, temporary holds, and conversions to open surgery remain part of the treatment-strategy estimand; withdrawal before surgery contributes available safety data and is explicitly classified in feasibility summaries; intraoperative discovery of unresectable disease is analyzed separately as a clinical-pathway outcome rather than silently excluded; and rescue interventions are retained within safety summaries because they are part of real-world risk management. Secondary endpoints are summarized descriptively with exact confidence intervals, and exploratory signals are interpreted as hypothesis-generating only.
}
\GrantTextArea[
  id=ops-sample-size, font=10pt, height=2.4in, maxchars=1550
]{Sample-size justification, assumptions, power or precision, and simulation strategy}{The planned sample size of 13 participants is justified as an early-phase safety-and-feasibility cohort, not a formal efficacy-powered comparison \citep{phase1ind}. This size is operationally consistent with a first-in-human program requiring staged review and sufficient exposure to estimate whether unacceptable toxicity, major procedural failure, or unmanageable cyber-physical risk are occurring at a rate that would preclude further development. Under exact binomial logic, 0 composite primary-endpoint events among^^J13 participants would place the upper bound of the two-sided 95\% confidence interval at approximately 24.7\%; 1 event would correspond to an upper bound of approximately 36.0\%, which is still informative for go/no-go decisions in a bounded Phase 1 setting. Design assumptions are therefore framed in terms of precision and decision utility rather than conventional power. Simulation work from the in silico, QSP, and digital-twin PDAC portfolio will be used prospectively to stress-test accrual assumptions, visit timing, missingness, replacement rules, perioperative scheduling, and stopping thresholds, but not to claim^^Jclinical efficacy or substitute for observed patient data \citep{qspmetpancre,chatgpt100kp,pdacdigtwinp,fdadigtwinpc}.
}
\GrantTextArea[
  id=ops-randomization, font=10pt, height=1.8in, maxchars=1450
]{Randomization, allocation concealment, masking, controls, and unblinding}{No randomization is used because this is a first-in-human, single-arm Phase 1 study intended to estimate safety, tolerability, and feasibility rather than comparative efficacy \citep{phase1ind,onpremwhippl}. Accordingly, allocation concealment is not applicable, the study is open label, and formal participant or investigator masking is not practical for the combined perioperative drug and robotic procedure. There is no concurrent control arm; external benchmarks and the later randomized Phase 2 protocol inform interpretation but do not function as statistical controls in this study \citep{daraxwhipple}. To reduce bias despite the absence of blinding, key operational and safety triggers will be prespecified, adverse-event coding will follow standard terminology, and independent review by the medical monitor or safety committee will be used for dose-escalation or continuation decisions. Unblinding rules are not applicable beyond restricted access to protected source records.
}
\GrantTextArea[
  id=ops-analysis-populations, font=10pt, height=2.25in, maxchars=1700
]{Primary and secondary analysis populations, missing data, multiplicity, and sensitivity analyses}{The primary safety population will include all enrolled participants who receive any study daraxonrasib exposure. A procedure population will include participants who proceed to the intended operative encounter, and a treated-protocol or DLT-evaluable population will include participants with sufficient exposure to contribute meaningfully to the primary assessment window unless an early event itself^^Jdefines evaluability. Secondary surgical and biomarker analyses will use endpoint-specific evaluable sets that are defined prospectively in the SAP \citep{phase1ind,onpremwhippl}. Missing data will be minimized through near-real-time source review, remote follow-up options, and prespecified window management; safety endpoints will prioritize conservative ascertainment, while descriptive secondary endpoints will report missingness transparently rather than relying on unjustified imputation. No formal multiplicity adjustment is planned for secondary or exploratory endpoints because they are descriptive^^Jand hypothesis-generating. Sensitivity analyses will examine the impact of replacements, open conversion cases, early withdrawal, and alternative endpoint-window definitions.
}
\GrantTextArea[
  id=ops-recruitment, font=10pt, height=2.4in, maxchars=1800
]{Recruitment, outreach, enrollment, retention, missed visits, and loss-to-follow-up mitigation}{Recruitment will focus on a high-volume pancreatic oncology catchment through multidisciplinary tumor board review, surgical oncology and medical oncology referral pathways, molecular prescreening for KRAS G12 status, and rapid pre-eligibility confirmation prior to consent \citep{phase1ind,onpremwhippl}. Outreach materials will be plain-language, culturally attentive, and available in translated form when needed. Enrollment feasibility will be supported by prespecified screening logs, weekly startup huddles during active recruitment, and replacement rules for screen failures or pre-treatment withdrawal. Retention will be strengthened through participant navigation, visit coordination around standard-of-care appointments, reminder systems, travel or lodging support when feasible, and remote or local-lab follow-up options for appropriate post-discharge assessments. Missed visits will trigger same-week contact attempts, windowed rescheduling, investigator review of clinically critical assessments, and use^^Jof routine care data when permitted. Loss to follow-up mitigation includes multiple contact methods, designated backup contacts, proactive scheduling of postoperative visits before discharge, and rapid outreach after any missed safety assessment.
}
\GrantTextArea[
  id=ops-safety, font=10pt, height=2.25in, maxchars=1800
]{Data and safety monitoring plan; DSMB / medical monitor; stopping and escalation rules}{Because this is a small, high-risk early-phase interventional study, the core oversight structure will consist of the sponsor-investigator, an independent medical monitor, and a standing safety review committee that functions as a DSMB-equivalent for cohort progression decisions \citep{phase1ind,onpremwhippl}. Safety data will be reviewed in real time during hospitalization and at prespecified intervals thereafter, with expedited review of serious adverse events, suspected unexpected serious adverse reactions, major surgical complications, device malfunctions, and cybersecurity or verification gate failures. Escalation or continuation to the next cohort or activation stage occurs only after documented review of DLTs, cumulative AEs, surgical recovery, device performance, and protocol deviations. Enrollment will pause for any treatment-related death, unexpected cluster of severe toxicity, uncontrolled major procedure-related harm, failure of essential safety overrides, or concern raised by the medical monitor, IRB, or FDA. Stopping rules will be finalized in the SAP and DSM plan, with clear authority for temporary hold, protocol modification, de-escalation, or termination.
}
\GrantTextArea[
  id=ops-quality, font=10pt, height=2.15in, maxchars=1550
]{Site activation, training, quality management, monitoring, auditing, and corrective actions}{Site activation will occur only after regulatory clearance, IRB approval, final SOP approval, credentialing of all key personnel, successful pharmacy and device accountability setup, dry-run simulations, and verification of source-to-eCRF data flow \citep{phase1ind,onpremwhippl,paisitedocpk}. Training will cover protocol procedures, GCP, consent, perioperative workflow, adverse-event grading, software^^Jand version control, emergency override and downtime procedures, cybersecurity hygiene, documentation standards, and protocol-deviation reporting. Quality management will use a risk-based^^Jplan that prioritizes eligibility confirmation, drug accountability, operative execution, endpoint adjudication, and serious-adverse-event reporting. Monitoring will include startup qualification, targeted remote and on-site review, central data checks, audit-trail review, and reconciliation across source records, EDC, pharmacy, and device logs. Any major finding will trigger documented corrective and preventive action, retraining, root-cause analysis, and re-verification before unrestricted continuation.
}
\GrantTextArea[
  id=ops-timeline, font=10pt, height=2.15in, maxchars=1550
]{Operational timeline, enrollment rate, milestones, contingencies, and decision gates}{Months 0-6 will be used for regulatory submission/clearance, IRB approval, PRS setup, final contracting, site qualification, simulation-based readiness testing, and activation. Months 7-18 will target enrollment of 13 participants, approximately 0.8 to 1.0 participants per month once screening is fully active, with cohort-level safety reviews after first 3, 6, 9, and 13 evaluable participants \citep{phase1ind,onpremwhippl}. Months 19-30 cover completion of the primary 30-day endpoint window, 90-day postoperative summaries, data cleaning, and database lock. Decision gates include FDA/IRB clearance before first participant, successful completion of activation drills before first procedure, acceptable cumulative safety before cohort progression, and an integrated go/no-go review after the full Phase 1 set. Contingencies include backup referral pipelines for slow accrual, preplanned protocol clarifications for avoidable deviations, replacement of failed hardware components, and temporary enrollment holds if safety, data quality, or regulatory concerns emerge.
}
